% =====================================================================
% Section II -- Background and Related Work
% SpecAsync-UVM manuscript, draft 1
%
% -------------------------------------------------------------------
% CITATION KEY RECONCILIATION -- READ BEFORE COMPILING
% -------------------------------------------------------------------
% BIB_AUDIT.md identified six entries needing repair, and T4_REPORT.md
% Section 1.2 records that all six were applied to ref.bib.  Whether the
% applied fix RENAMED the keys or upgraded them in place is not recorded,
% so the keys used below are a best reconstruction.  Check these six and
% sed as needed before the first build:
%
%   OLD KEY                          KEY USED HERE
%   anonymous2022learningoversub  -> long2022oversub          (replaced entry)
%   jain2019crac                  -> jain2020crac             (replaced entry)
%   bastemTilesTR                 -> bastem2017tiles          (replaced entry)
%   long2022dlprefetch            -> long2022dlprefetch       (in-place upgrade to JPDC 2023)
%   choi2022gpuaware              -> choi2022gpuaware         (in-place upgrade to IPDPSW 2022)
%   parravicini2021dagpolyglot    -> parravicini2021dagpolyglot (in-place upgrade to IPDPS 2021)
%
% -------------------------------------------------------------------
% EVALUATION-SUBSTRATE COLUMN -- VERIFY BEFORE SUBMISSION
% -------------------------------------------------------------------
% The table's "Evaluated on" column is the load-bearing new element and
% directly answers Reviewer Comment 5.  Entries marked below are asserted
% from each paper's own description as summarised in the previous draft;
% the following four were NOT independently re-verified and should be
% checked against the papers themselves before submission:
%   CONFIRMED hardware: go2023earlyadaptor (abstract states real hardware,
%     NVIDIA RTX 3090, open-source driver), shen2026stride (non-invasive
%     runtime layer, no kernel/application modification), guo2025ghost,
%     jung2020hum, allen2021indepth, jones2025helm, wahlgren2025mi300a.
%   CONFIRMED simulation: lin2025forest, long2022dlprefetch, yu2020cppe.
%   STILL UNVERIFIED -- check before submission:
%     li2019oversubframework (ASPLOS'19), pratheek2024suv (MICRO'24),
%     nihaal2024compression (ICPP'24).
% Getting this column wrong would be a serious and easily-caught error, so
% it is better to drop a row than to guess at its substrate.
% =====================================================================

\section{Background and Related Work}
\label{sec:related}

Work on reducing UVM fault-handling cost falls into four groups: measurement
studies that characterise where fault-service time goes, prediction and
prefetching schemes that try to move data earlier, oversubscription policies
that manage residency under memory pressure, and asynchronous-overlap
techniques that hide transfer cost behind computation. Our study sits at the
intersection of the first and the last: it builds an overlap mechanism and then
measures why it fails. We therefore organise this section around a question that
cuts across all four groups---\emph{what was each result measured on}---because
that question turns out to be what separates our finding from the expectations
the literature sets.

\subsection{Where Fault-Service Time Goes}

The closest relatives to this work are direct measurements of the UVM fault
path. Allen and Ge first characterised UVM at the driver level, establishing
the fault batch as the fundamental unit of work and showing that batch service
time dominates migration cost \cite{allen2021indepth}. Their subsequent work
provides a fine-grained decomposition of demand paging in
UVM, quantifying per-batch phases and showing that the \emph{service} portion---%
host-side unmapping and page-table work---can dominate CPU--GPU migration cost,
while DMA is often overlapped in the background \cite{allen2024finegrain}. Their
analysis also surfaces a trade-off we encounter directly: batching reduces trap
frequency, but larger batches amplify synchronization and tracker-wait effects,
capping end-to-end gains \cite{allen2024finegrain}. Our seven-phase
decomposition is methodologically in this tradition and agrees with its
qualitative conclusion, but differs in instrumentation point and in what it is
used for: we measure inside a modified driver rather than around it, and we use
the decomposition to derive an \emph{upper bound} on what an off-path scheme
could recover rather than to characterise the path descriptively.

The same group's more recent work derives locality metrics directly from UVM
driver page-fault telemetry to guide allocation under oversubscription
\cite{jones2025helm}. That study and ours draw on the same instrumentation
source and reach opposite kinds of conclusion---theirs constructive, ours
negative---which we take as complementary rather than conflicting: telemetry
that is informative enough to guide residency decisions is not thereby
sufficient to make off-path speculation profitable. Characterisation has also
begun to extend to integrated CPU--GPU parts, where page-fault handling and
TLB management differ substantially from the discrete-GPU path we study
\cite{wahlgren2025mi300a}; we return to what this bounds about our claims in
the Discussion.

Benchmark-level characterisation supplies the other half of the picture.
UVMBench contributes a suite spanning machine learning, graph, HPC, and
statistical workloads, and shows that naive UVM can underperform non-UVM code
while explicit prefetching largely closes the gap \cite{gu2020uvmbench}. Chien
\emph{et al.} evaluate advanced Unified Memory features and find that
advice and prefetch choices materially affect performance, but in
platform-dependent ways, with oversubscription regimes degrading results
\cite{chien2019umfeatures}. Both findings recur in our data. The
platform-dependence in particular is why we report every cost ratio with an
explicit platform qualifier rather than as a property of UVM in general.

A separate line moves the paging control plane itself. GPUVM shifts portions of
unified-memory control toward the GPU using RDMA-assisted mechanisms, reporting
substantial latency improvements on graph workloads
\cite{nazaraliyev2024gpuvm}. This is a more radical restructuring than ours and
depends on specialised networking assumptions, but it shares our conclusion in
one respect: it treats the host-side path as the thing to be removed rather than
the thing to be predicted around.

\subsection{Prediction and Prefetching}

A large body of work improves UVM performance by initiating movement earlier.
Early-Adaptor monitors fault behaviour and adapts proactive UVM actions
dynamically, reporting speedups under oversubscription by adjusting prefetch
aggressiveness \cite{go2023earlyadaptor}. SUV uses static analysis to guide UVM
management---selectively pinning dense data and issuing iteration-aligned
prefetch---achieving large speedups over baseline UVM
\cite{pratheek2024suv}. Both indicate that the \emph{timing} of proactive
actions matters, not merely their presence.

Learning-based predictors push this further. Long \emph{et al.} propose a deep
learning prefetcher built from transformer-style components, improving IPC and
page hit rate \cite{long2022dlprefetch}; Forest configures tree prefetchers at
runtime according to each data object's access pattern, on the observation that
the stock access-pattern-oblivious tree-based neighbouring prefetcher causes
excessive thrashing and unnecessary migration \cite{lin2025forest}. That
observation is worth dwelling on, because it concerns the same in-path
prefetcher whose interaction with our mechanism we report in
Section~\ref{sec:squeeze}: the component is simultaneously the reason
speculation has nothing to do when it is enabled and, by Forest's account, an
imperfect policy in its own right.

Most directly relevant to our stride policy, Shen \emph{et al.} report a
stride-aware Markov prefetcher for GPU unified memory that reaches high
prediction accuracy and delivers real speedups, prototyped as a non-invasive
runtime layer on real workloads \cite{shen2026stride}. Our stride policy, by
contrast, measurably \emph{harms} performance. The two results are compatible,
and the reason they are compatible is the point of this paper. Their predictor
drives migration from above the driver and never pays a per-fault dispatch
cost; ours issues speculative work from inside the fault handler and pays that
cost on every fault. More fundamentally, our negative result is not a statement
about prediction accuracy at all---our oracle achieves 100\% accuracy by
construction and still yields nothing. A predictor can therefore be as accurate
as Shen \emph{et al.} report and still be unable to help when its output is
consumed off-path inside the driver.

Our oracle policy can be read as the limiting case of this whole line: it makes
prediction accuracy perfect by construction and so bounds from above what any
predictor in this family can deliver on this driver. That the bound is empty is
the result this paper reports.

\subsection{Oversubscription and Memory Pressure}

Under oversubscription, UVM enters eviction-dominated regimes where performance
collapses. Li \emph{et al.} propose coordinated admission and throttling
policies to reduce thrashing \cite{li2019oversubframework}, and Yu \emph{et al.}
couple prefetch with modified eviction, reporting large speedups
\cite{yu2020cppe}. Others reduce effective pressure by increasing usable
capacity through selective compression and packing \cite{nihaal2024compression},
or through oversubscription-aware scheduling of tensor migration
\cite{kim2025most}. Long \emph{et al.} approach the same regime with a learned
oversubscription-management framework \cite{long2022oversub}.

Our oversubscription measurements are complementary rather than competitive: we
do not propose a residency policy, we measure what happens when speculation is
run inside one. The result is instructive. Under thrashing, speculation behaves
as an extremely costly accidental page-replacement policy---in one
configuration issuing over 80 million speculative migrations to service fewer
than two thousand recorded hits---rather than as working prediction. This is
precisely the regime in which the coordinated policies above are valuable, and
precisely the regime in which an uncoordinated off-path mechanism is most
harmful.

\subsection{Asynchronous Overlap and Off-Path Work}

Finally, a body of work shows that explicit overlap of transfer and computation
delivers real gains, generally above the driver. HUM converts synchronous
host-to-device copies into an asynchronous pipeline using page-fault-driven
mechanisms \cite{jung2020hum}, and tile-based partitioning overlaps transfers
with computation using caching and ghost regions \cite{bastem2017tiles}.
Runtime systems generalise this: HPXCL integrates CUDA execution into an
asynchronous tasking framework with futures \cite{diehl2018hpxcl}; GPU-aware
asynchronous tasks with prioritised streams and kernel fusion improve
scalability by increasing overlap \cite{choi2022gpuaware}; DAG-based scheduling
in a polyglot runtime overlaps transfers using runtime-inferred dependencies
\cite{parravicini2021dagpolyglot}. CRAC, while targeting checkpoint-restart
rather than performance, documents the engineering hazards of coordinating
around CUDA streams and UVM \cite{jain2020crac}. Outside GPUs entirely,
pre-scheduling has been used to cut completion-time variance by moving
coordination work earlier \cite{jin2019prescheduling}.

Two recent results sharpen the comparison. Ghost Threading demonstrates
software-only helper-thread prefetching on real systems, and identifies
\emph{timeliness} as the governing constraint: the helper thread must run far
enough ahead to be useful but not so far as to be wasteful, and achieving that
required a synchronisation primitive cheap enough to throttle the helper thread
without OS involvement \cite{guo2025ghost}. AGIO makes an analogous argument
for GPU-side I/O, decoupling initiation from completion so that microsecond-scale
SSD latencies can be overlapped with useful work \cite{han2026agio}.

We regard these as corroboration rather than counterexample, and they are the
clearest statement of the condition our mechanism fails to meet. Off-path work
pays when the cost of dispatching it is small relative to the latency being
hidden. Ghost Threading found a primitive cheap enough at instruction
granularity; AGIO has whole SSD accesses to hide behind. Inside the UVM fault
handler neither holds: the dispatch cost and the work being overlapped are the
same order of magnitude, and under load the dispatch window inflates by three
orders of magnitude while faults continue to arrive. Our contribution is to
show that the principle does not survive being applied at the granularity of an
individual page fault inside a production driver, and to measure the margin by
which it fails.

\subsection{Limit Studies and Measurement Methodology}

Two methodological literatures underpin this study and are worth naming
explicitly, because both of our less conventional choices---an oracle policy
and a catalog of our own measurement errors---have established precedent.

Bounding achievable benefit by granting a mechanism oracular knowledge of
future accesses is a recognised technique. R-Max applies exactly this idea to
CPU cache prefetching, approximating an ideal prefetcher under realistic
bandwidth and capacity constraints in order to quantify how much gain remains
unexploited by current designs \cite{wang2025rmax}. Our oracle serves the same
function one level down the memory hierarchy. The contrast in outcome is
instructive: that work finds substantial headroom left on the table for CPU
caches, whereas we find the UVM fault path affords between 0.08\% and 7.94\%.
The method is the same; the bound is a property of the subsystem, not of the
technique.

The measurement-pitfalls catalog we present in Section~\ref{sec:pitfalls} also
has precedent. Van der Kouwe \emph{et al.} identify twenty-two benchmarking
flaws and, surveying fifty defence papers at top venues, find that tier-one
papers contain an average of five flaws each \cite{vanderkouwe2019sok}.
Papadopoulos \emph{et al.} propose eight principles for reproducible
performance evaluation and find, by systematic review, that few published
studies follow them \cite{papadopoulos2019principles}. Raasveldt \emph{et al.}
demonstrate with mock benchmarks how easily large apparent differences can be
produced where none exist \cite{raasveldt2018fair}---structurally the same
hazard as the entry in our catalog where the measurement and the arithmetic
were both correct and the conclusion was still wrong, because the comparison
was made against a crippled baseline. Our pre-registration, strictly
interleaved run ordering, and reporting of minimum detectable effects are
responses to this literature rather than local idiosyncrasies.

\subsection{The Simulation--Hardware Gap}

Table~\ref{tab:prior-work} adds a column that the previous version of this
survey did not carry: what each result was evaluated on. The pattern it exposes
is the motivation for this paper. Measurement studies
\cite{allen2024finegrain,gu2020uvmbench,chien2019umfeatures} run on real
hardware and report costs. Mechanism proposals that predict, prefetch, or manage
residency are substantially more likely to be evaluated in simulation, where the
proposed mechanism's cost model is written by the same author who wrote the
mechanism.

We want to be careful about what this observation does and does not license.
Simulation is not a methodological failing; it is often the only tractable way
to explore a design space, and several of the simulated results above are
carefully done. The point is narrower and specific to our mechanism: off-path
speculation's benefit case rests entirely on a comparison between two costs---%
the latency of dispatching speculative work and the latency of the fault it is
meant to hide. A simulator that models fault-service time but abstracts the
driver's internal queueing, its existing in-path prefetcher, and its
address-space lock discipline will get that comparison wrong in a direction that
favours the mechanism. All three of those components turned out to be decisive
in our measurements: the in-path prefetcher makes speculation redundant when
enabled and unaffordable to disable, the queueing inflates dispatch latency by
three orders of magnitude under load, and the lock-held phase is CPU-side
bookkeeping rather than a GPU wait that asynchrony could hide.

This is the gap our study addresses. Not that prior work is wrong---the
prediction results in particular we take as sound, and our oracle presupposes
them---but that a mechanism whose viability depends on driver-internal cost
structure has to be measured against that structure before it can be believed.

% ---------------------------------------------------------------------
% Prior-work table.  The "Evaluated on" column is the new element; see the
% verification note in this file's header before submission.
% ---------------------------------------------------------------------
\begin{table*}[t]
\centering
\footnotesize
\setlength{\tabcolsep}{4pt}
\renewcommand{\arraystretch}{1.15}
\newcolumntype{Y}{>{\raggedright\arraybackslash}X}

\begin{tabularx}{\textwidth}{@{}l Y Y l Y@{}}
\toprule
\textbf{Work} &
\textbf{Focus} &
\textbf{Mechanism} &
\makecell[l]{\textbf{Evaluated}\\\textbf{on}} &
\makecell[l]{\textbf{Relation to}\\\textbf{this study}} \\
\midrule

\cite{allen2024finegrain} &
UVM fault cost breakdown &
Measurement, driver-level &
Hardware &
Closest prior measurement; we decompose inside a modified driver and derive a bound \\

\cite{gu2020uvmbench} &
UVM benchmarking &
Benchmark suite &
Hardware &
Establishes prefetch closes the naive-UVM gap; consistent with our squeeze finding \\

\cite{chien2019umfeatures} &
Unified Memory features &
Measurement &
Hardware &
Platform-dependence motivates our per-platform qualification of every cost ratio \\

\cite{nazaraliyev2024gpuvm} &
Alternative paging control plane &
RDMA-assisted GPU-side control &
Hardware &
Also treats the host path as the target; assumes specialised networking \\

\cite{go2023earlyadaptor} &
Proactive UVM management &
Adaptive runtime policy &
Hardware &
Timing of proactive action matters; complements service-path reduction \\

\cite{pratheek2024suv} &
Static-analysis-guided UVM &
Compiler analysis + policy &
Hardware &
Program-informed prefetch; a candidate predictor input, bounded by our oracle \\

\cite{long2022dlprefetch} &
Learned prefetching &
Deep learning predictor &
Simulation &
Prediction is learnable; our oracle upper-bounds what any such predictor can yield here \\

\cite{lin2025forest} &
Access-aware management &
Prediction-driven queuing &
Simulation &
Same bound applies; interaction with driver queueing not modelled \\

\cite{yu2020cppe} &
Coordinated prefetch/eviction &
Policy co-design &
Simulation &
Residency coordination, not host-service cost \\

\cite{nihaal2024compression} &
Oversubscription capacity &
Selective compression &
Simulation &
Reduces pressure rather than service cost; orthogonal \\

\cite{jones2025helm} &
UM access characterisation &
Driver-telemetry metrics + policy &
Hardware &
Same telemetry source as this work; constructive result where ours is negative \\

\cite{shen2026stride} &
UVM page prefetching &
Stride Markov predictor, runtime layer &
Hardware &
Accurate stride prediction above the driver; our stride policy harms below it \\

\cite{guo2025ghost} &
Helper-thread prefetching &
Software-only, SMT helper threads &
Hardware &
Identifies timeliness and cheap throttling as the enabling condition we lack \\

\cite{jung2020hum} &
Transfer overlap &
Fault-driven async pipeline &
Hardware &
Strongest support for off-path principle; overlaps at transfer, not per-fault, granularity \\

\cite{bastem2017tiles} &
Transfer/compute overlap &
Tiling + ghost regions &
Hardware &
Same principle at tile granularity \\

\bottomrule
\end{tabularx}
\caption{Representative prior work, classified by evaluation substrate.
Mechanism proposals evaluated in simulation are systematically more optimistic
about off-path benefit than hardware measurement supports, because the costs
that defeat the mechanism---in-path prefetching, driver queueing, and
lock-held CPU bookkeeping---are typically not modelled.}
\label{tab:prior-work}
\end{table*}
